\documentclass[11pt,a4paper]{article}

%─── Packages ─────────────────────────────────────────────────────────────────
\usepackage{amsmath,amssymb,amsthm}
\usepackage{geometry}
\usepackage{booktabs}
\usepackage[bookmarks=false]{hyperref}
\usepackage{array}

\usepackage{xcolor}
\usepackage{enumitem}
\usepackage{graphicx}
\usepackage{float}
\geometry{top=2.5cm,bottom=2.8cm,left=2.5cm,right=2.5cm}

%─── Theorem environments ─────────────────────────────────────────────────────
\newtheorem{theorem}{Theorem}[section]
\newtheorem{lemma}[theorem]{Lemma}
\newtheorem{proposition}[theorem]{Proposition}
\newtheorem{corollary}[theorem]{Corollary}
\theoremstyle{definition}
\newtheorem{definition}[theorem]{Definition}
\theoremstyle{remark}
\newtheorem{remark}[theorem]{Remark}
\newtheorem{conjecture}{Conjecture}

%─── Macros ───────────────────────────────────────────────────────────────────
\newcommand{\vp}{\varphi}
\newcommand{\lam}{\lambda}
\newcommand{\bst}{\beta^{*}}
\newcommand{\sinW}{\sin^{2}\!\theta_{W}}
\newcommand{\TCMB}{T_{\mathrm{CMB}}}
\newcommand{\vEW}{v_{\mathrm{EW}}}
\newcommand{\Lcond}{\Lambda_{\mathrm{cond}}}
\newcommand{\LUV}{\Lambda_{\mathrm{UV}}}
\newcommand{\MPl}{M_{\mathrm{Pl}}}
\newcommand{\LQCD}{\Lambda_{\mathrm{QCD}}}
\newcommand{\ERy}{E_{\mathrm{Ry}}}
\newcommand{\aem}{\alpha_{\mathrm{em}}}
\newcommand{\eps}{\varepsilon}
% New macros for this paper
\newcommand{\Hcond}{H_{\mathrm{cond}}}
\newcommand{\Heff}{H_{\mathrm{eff}}}
\newcommand{\rhoa}{\rho_{\mathrm{atom}}}
\newcommand{\rhoinf}{\rho_{\infty}}
\newcommand{\Seff}{S_{\mathrm{eff}}}
\newcommand{\epsr}{\varepsilon(r)}

%─── Colours ──────────────────────────────────────────────────────────────────
\definecolor{darkgreen}{RGB}{0,120,60}
\definecolor{darkred}{RGB}{160,20,20}
\definecolor{darkblue}{RGB}{0,40,140}
\definecolor{darkorange}{RGB}{180,80,0}

%══════════════════════════════════════════════════════════════════════════════
\title{\textbf{Atomic Quantum Mechanics from the Hopfion Condensate}
\\[0.6em]
{\large Companion Paper XIII to the Density-Feedback
Faddeev--Niemi Hopfion Series}
\\[0.6em]
\large Preprint --- June 2026}
\author{Frederick Manfredi}
\date{}
%══════════════════════════════════════════════════════════════════════════════

\begin{document}
\maketitle

\paragraph*{Units and conventions.}
Natural units $\hbar = c = k_B = 1$ throughout.
Energies in eV unless stated otherwise.
$\MPl = 1.22\times10^{19}\,\mathrm{GeV}$ (unreduced Planck mass).
The condensate scale is
$\Lcond = \TCMB(\pi^2/150)^{1/4} = 1.1895\times10^{-4}\,\mathrm{eV}$.
The fine-structure constant and electron mass are framework outputs:
$\aem^{-1} = 137.035\,998\,486$ (Paper~IV~\cite{Paper4}, four-term,
$5\times10^{-7}\%$) and
$m_e = 0.51100\,\mathrm{MeV}$
(Paper~IV~\cite{Paper4}, Paper~XII~\cite{Paper12}).
The condensate feedback coupling is $\bst\approx0.452$
(Paper~I~\cite{Paper1}, self-consistent saddle point).

%─── Abstract ─────────────────────────────────────────────────────────────────
\begin{abstract}
We derive the condensate-modified Schr\"odinger equation for an electron
gradient-knot in the density-feedback Faddeev--Niemi Hopfion background
in the presence of a nuclear Coulomb source.
The suppression law $\Seff(\theta,\rho)=\sin^4\!\theta/[\vp^6(1+\bst\rho)]$
of Paper~I adds a single correction to the standard hydrogen Hamiltonian:
\[
  \Heff \;=\; -\frac{\hbar^2}{2m_e}\nabla^2 - \frac{e^2}{r}
  \;-\; \frac{\hbar^2}{2m_e}\,\frac{\epsr}{r^2}\,L^2,
  \qquad
  \epsr = \frac{1}{\vp^6(1+\bst\rho(r))},
\]
where the condensate correction applies only to angular kinetic energy
(the radial correction is identically zero by $\Seff(0,\rho)=0$).

Two regimes determine the physics.
\textbf{Inside atoms}: the atom's own electromagnetic field drives
$\bst\rho_{\mathrm{atom}}(a_0)\sim\aem^6 m_e^4/(2\Lcond^4)\approx2.6\times10^{25}$,
chameleon-screening the condensate correction to
$\eps(a_0)\approx2.2\times10^{-27}$.
The orbital shapes are standard spherical harmonics
$|Y_\ell^m(\theta,\varphi)|^2$ to any measurable precision, and the
condensate correction to hydrogen energy levels is $\sim10^{-24}\,\mathrm{eV}$
(suppressed $\sim10^{20}$ below the fine structure).
The spherical harmonics of atomic physics are thus \emph{derived} within
the condensate framework, not assumed: they emerge from the chameleon
mechanism that also resolves the cosmological constant hierarchy.

\textbf{Free electrons in the cosmic background}: with $\bst\rho_\infty = \vp$
(exact at the condensate attractor), $(1+\bst\rho_\infty) = \vp^2$
(exact, from $\vp^2=\vp+1$), and $\eps = 1/\vp^8$.
The condensate is anisotropic at the $1/\vp^8\approx2.13\%$ level.
The effective mass for tangential propagation is
$m_t = m_e(1+1/\vp^8) = m_e(\vp^8+1)/\vp^8$
while radial propagation has $m_r = m_e$ exactly.
This $2.13\%$ mass anisotropy is a falsifiable prediction for
free-electron propagation in low-density environments
(e.g.\ precision Penning-trap measurements oriented relative to the
local condensate gradient or the CMB dipole).

The same $(1+\bst\rho)$ denominator screens both effects: it drives
the cosmological constant below the observed bound at cosmic density
and forces orbital shapes to be spherical harmonics at atomic density.
The chameleon screening extends universally to all $Z$-electron atoms:
the condensate correction is bounded by $\sim10^{-24}\,\mathrm{eV}$
for all atoms, scaling as $Z^{-4}$ for inner electrons
and leaving the Hartree--Fock equations exact.

Finally, the fermionic nature of the electron is derived from the
condensate topology.
The binary icosahedral group $2I$, the symmetry group of the condensate
(Paper~I), is a double cover of the icosahedral group $I$.
The $Q=2$ gradient-knot transforms in the fundamental spinor representation
of $2I$: the central element $z$ (the $2\pi$ rotation) acts as $-1$,
giving spin $J=\tfrac{1}{2}$ (Theorem~\ref{P13:thm:spin_half}),
Fermi statistics (Corollary~\ref{P13:cor:fermi_stats}), and
Pauli exclusion (Corollary~\ref{P13:cor:pauli_exclusion}).
This is the same $\mathbb{Z}_2$ structure that underlies the
colour level shift $\delta n^{(C)}=\tfrac{1}{2}$ of Paper~V.\end{abstract}

\footnotesize
\noindent\textbf{Keywords:} condensate-modified Schr\"odinger equation;
chameleon screening; spherical harmonics; hydrogen atom; multi-electron atoms;
Hartree--Fock; free-electron mass anisotropy; density-feedback Hopfion;
effective Hamiltonian; orbital shapes; Penning trap;
spin-$\tfrac{1}{2}$; Fermi statistics; Pauli exclusion principle;
binary icosahedral group $2I$; Hopf charge; spin tower $J=Q/4$;
Born--Oppenheimer approximation; non-adiabatic coupling;
translation fibration; WZW model; Standard Model spin assignments.
\normalsize

\tableofcontents
\bigskip

\newpage
%══════════════════════════════════════════════════════════════════════════════
\section{Introduction}
\label{P13:sec:intro}
%══════════════════════════════════════════════════════════════════════════════

Papers~I--XIV of this series establish that a single topological object
--- the density-feedback Faddeev--Niemi Hopfion condensate --- derives
approximately 45 Standard Model, gravitational, and cosmological
observables from a single experimental input $\TCMB$
(Paper~XII~\cite{Paper12}, one-parameter chain).
The suppression law
\begin{equation}
  \Seff(\theta,\rho)
  \;=\; \frac{\sin^4\!\theta}{\vp^6(1+\bst\rho)},
  \label{P13:eq:suppression}
\end{equation}
governs how field-gradient energy propagates at angle $\theta$ to
the local preferred direction in a medium of condensate density $\rho$.
The Bogomolny parameter $\vp^6$ is fixed by the icosahedral topology
($Q_{\rm num}=Q_{\rm group}=10$, Paper~I~\cite{Paper1});
the feedback coupling $\bst\approx0.452$ is the unique
self-consistent saddle-point value.

The same $(1+\bst\rho)$ denominator appears throughout:
it screens the condensate's vacuum energy below the observed
cosmological constant bound (Paper~VII~\cite{Paper7}, chameleon mechanism)
and it suppresses the condensate's coupling to high-energy processes
in QCD and electroweak environments (Papers~IV--V~\cite{Paper4,Paper5}).
In each case the physical result is protected from condensate
corrections by the fact that the local density far exceeds the
cosmic attractor $\rho_\infty=\vp/\bst$.

This paper asks: what does the condensate do to atomic physics?
The electron is a gradient-knot whose mass is fixed by the condensate
(Papers~I, IV).
When orbiting a proton, the condensate is chameleon-screened to
$\eps_{\rm cond}\sim10^{-27}$, leaving standard quantum mechanics exact.
Far from matter, in the cosmic condensate background at $\bst\rho_\infty=\vp$,
the suppression lifts and the condensate is anisotropic at the 2.13\% level.

The paper establishes ten results, all without additional postulates:
\begin{enumerate}[noitemsep]
\item The condensate-modified Hamiltonian (Theorem~\ref{P13:thm:modified_H}).
\item Spherical harmonics as exact orbital shapes from chameleon screening
  (Theorem~\ref{P13:thm:chameleon_screening}).
\item Condensate energy corrections $\Delta E\lesssim10^{-24}\,\mathrm{eV}$
  (Proposition~\ref{P13:prop:energy_correction}).
\item Free-electron mass anisotropy $\Delta m/m_e=1/\vp^8\approx2.13\%$,
  a falsifiable Penning-trap prediction (Proposition~\ref{P13:prop:mass_anisotropy}).
\item Universality for all $Z$-electron atoms: spherical harmonics hold,
  $\Delta E$ bounded by $\sim10^{-24}\,\mathrm{eV}$, Hartree--Fock exact
  (Propositions~\ref{P13:prop:multi_electron_screening},
  \ref{P13:prop:multi_electron_energy}).
\item The Schr\"odinger equation from condensate quantisation, exact to
  $O(\eps_{\rm cond})\sim10^{-27}$ (Theorem~\ref{P13:thm:schrodinger}).
\item Exact vanishing of Born--Oppenheimer non-adiabatic coupling
  ($d_{0k}=0$ for all $k$, by the translation fibration), tightening
  the error bound by $\sim10^{17}$
  (Theorem~\ref{P13:thm:NA_vanishing}, Corollary~\ref{P13:cor:improved_BO}).
\item Spin-$\tfrac{1}{2}$, Fermi statistics, and Pauli exclusion from
  $Q=2$ topology and the binary icosahedral group $2I$
  (Theorems~\ref{P13:thm:spin_half}--\ref{P13:cor:pauli_exclusion}).
\item Complete spin assignment $J=Q/4$ for all Standard Model particles
  from the $Q$-tower and WZW energy minimisation
  (Theorem~\ref{P13:thm:spin_tower}, Table~\ref{P13:tab:spin_assignments}).
\item A complete derivation of atomic and molecular quantum mechanics ---
  orbital shapes, the Schr\"odinger equation, Pauli exclusion, Hartree--Fock,
  quantum chemistry --- from the condensate framework
  (Theorem~\ref{P13:thm:atomic_foundation}).
\end{enumerate}

\paragraph{Position in the series.}
This paper depends on Papers~I--IV and~XII for the condensate parameters
and the electron mass; Paper~X for the condensate Hilbert space and
quantisation; Paper~V for the $2I$ character table, WZW level $k=3$,
and the colour level shift $\delta n^{(C)}=\tfrac{1}{2}$.

\paragraph{Organisation.}
Section~\ref{P13:sec:setup} sets up the electron gradient-knot in the
condensate background.
Section~\ref{P13:sec:modified_H} proves the modified Hamiltonian.
Section~\ref{P13:sec:chameleon} derives spherical harmonics from chameleon
screening and extends the result to all $Z$-electron atoms.
Section~\ref{P13:sec:energy_correction} computes energy-level corrections.
Section~\ref{P13:sec:anisotropy} derives the free-electron mass anisotropy.
Section~\ref{P13:sec:schrodinger} derives the Schr\"odinger equation and
proves the exact vanishing of non-adiabatic corrections.
Section~\ref{P13:sec:fermi} derives spin-$\tfrac{1}{2}$, Fermi statistics,
Pauli exclusion, and the spin assignment $J=Q/4$.
Section~\ref{P13:sec:open} states the remaining open problem.
Section~\ref{P13:sec:summary} summarises.

%══════════════════════════════════════════════════════════════════════════════
\section{Setup: The Electron Gradient-Knot in the Condensate}
\label{P13:sec:setup}
%══════════════════════════════════════════════════════════════════════════════

\subsection{The electron as a gradient-knot}

In the condensate framework (Paper~VII~\cite{Paper7}, Axiom~A3),
stable localised field configurations --- gradient-knots --- exist at
the discrete mass scales $m_n = m_0\vp^{-2n}$ of the $\vp$-tower.
The electron is the gradient-knot at level $n = 20$ (Papers~I and~IV):
$m_e = \Lcond\,\vp^{20}\,e^{4\pi-3/400-\aem/\pi}$
(Paper~IV~\cite{Paper4}, Theorem~8.11, $0.013\%$).

When the electron gradient-knot is localised near a proton, its
quantum state is described by a wavefunction $\psi(r,\theta,\varphi)$
giving the probability amplitude for the centre of the knot to be at
position $\mathbf{r}$.
The preferred direction at position $\mathbf{r}$ is
$\hat{r} = \mathbf{r}/|\mathbf{r}|$ (radially outward from the proton),
because the proton's Coulomb field --- and through it the condensate
density gradient --- is spherically symmetric and radially directed.

\subsection{The condensate density profile near the proton}

The condensate density $\rho(r)$ near the proton is determined by
local equilibrium between the condensate's self-interaction and the
electromagnetic field of the proton.
At distance $r$ from the proton, the electromagnetic energy density is
\begin{equation}
  u_{\rm EM}(r)
  \;=\; \frac{\aem^2}{2r^4}
  \;=\; \frac{\aem^6 m_e^4}{2}\cdot\left(\frac{a_0}{r}\right)^4,
  \label{P13:eq:uEM}
\end{equation}
in natural units ($\hbar=c=1$), where $a_0 = 1/(\aem m_e)$ is the
Bohr radius.
The condensate density is driven to equilibrium with this field
at the scale where the condensate and EM energy densities are comparable:
\begin{equation}
  \bst\rho_{\rm atom}(r)
  \;\sim\;
  \frac{u_{\rm EM}(r)}{\Lcond^4}
  \;=\; \frac{\aem^6 m_e^4}{2\Lcond^4}
        \left(\frac{a_0}{r}\right)^4.
  \label{P13:eq:rho_atom}
\end{equation}
At the Bohr radius ($r = a_0$):
\begin{equation}
  \bst\rho_{\rm atom}(a_0)
  \;\approx\; \frac{\aem^6 m_e^4}{2\Lcond^4}
  \;\approx\; 2.57\times10^{25},
  \label{P13:eq:rho_atom_a0}
\end{equation}
which far exceeds the cosmic attractor value
$\bst\rhoinf = \vp \approx 1.618$.
The atomic interior is in the \emph{chameleon high-density regime}
$\bst\rho\gg1$, the same regime that controls the cosmological
constant screening (Paper~VII~\cite{Paper7}).

\begin{remark}[Radial dependence]
\label{P13:rem:radial_dep}
Equation~\eqref{P13:eq:rho_atom} gives $\bst\rho_{\rm atom}(r)\propto r^{-4}$:
the condensate density grows rapidly toward the nucleus.
At $r = 5a_0$ (outer edge of the $n=2$ shell):
$\bst\rho_{\rm atom}\approx4.1\times10^{22}$, still $\gg1$.
The chameleon limit is an excellent approximation throughout the atom.
\end{remark}

%══════════════════════════════════════════════════════════════════════════════
\section{The Condensate-Modified Hydrogen Hamiltonian}
\label{P13:sec:modified_H}
%══════════════════════════════════════════════════════════════════════════════

\subsection{Derivation of the correction term}

The condensate's contribution to the kinetic energy of the electron
gradient-knot is determined by the suppression law~\eqref{P13:eq:suppression}.
For an electron wavefunction $\psi(\mathbf{r})$ with gradient
$\nabla\psi$, decompose the gradient into radial and angular components:
\begin{equation}
  |\nabla\psi|^2
  \;=\; \underbrace{\left|\frac{\partial\psi}{\partial r}\right|^2}_{\text{radial}}
  \;+\; \underbrace{\frac{1}{r^2}|L\psi|^2}_{\text{angular}},
  \label{P13:eq:gradient_decomp}
\end{equation}
where $L^2 = -r^2\nabla_\perp^2$ is the angular momentum operator
(eigenvalue $\ell(\ell+1)\hbar^2$ on $Y_\ell^m$).

The condensate energy for each component depends on the angle
$\theta_{\rm local}$ between the gradient and the preferred direction $\hat{r}$:
\begin{itemize}[noitemsep]
  \item \emph{Radial gradient}: $\theta_{\rm local}=0$,
    $\Seff(0,\rho) = \sin^4\!0/[\vp^6(1+\bst\rho)] = 0$ exactly.
  \item \emph{Angular gradient}: $\theta_{\rm local}=\pi/2$,
    $\Seff(\pi/2,\rho) = 1/[\vp^6(1+\bst\rho)] \equiv \epsr$.
\end{itemize}
The condensate correction to the kinetic energy of the electron
wavefunction is therefore
\begin{equation}
  \Hcond
  \;=\; -\frac{\hbar^2}{2m_e}
        \,\frac{\epsr}{r^2}\,L^2,
  \label{P13:eq:H_cond}
\end{equation}
where
\begin{equation}
  \boxed{\epsr
  \;\equiv\; \frac{1}{\vp^6(1+\bst\rho(r))}}
  \label{P13:eq:epsilon_def}
\end{equation}
is the local condensate suppression function.
Equation~\eqref{P13:eq:H_cond} is exact: the only approximation is the
identification of $\theta_{\rm local}$ with the spherical polar angle
relative to $\hat{r}$, which holds for a spherically symmetric source.

\begin{theorem}[Condensate-modified hydrogen Hamiltonian]
\label{P13:thm:modified_H}
The effective Hamiltonian for the electron gradient-knot in the
condensate background near a proton is
\begin{equation}
  \boxed{
  \Heff
  \;=\; -\frac{\hbar^2}{2m_e}\nabla^2 - \frac{e^2}{r}
  \;-\; \frac{\hbar^2}{2m_e}\,\frac{\epsr}{r^2}\,L^2,
  }
  \label{P13:eq:H_eff}
\end{equation}
where $\epsr = 1/[\vp^6(1+\bst\rho(r))]$ and $\rho(r)$ is the
condensate density at radius $r$.
The condensate correction vanishes identically for radial motion
($\Seff(0,\rho)=0$ by $\sin^4\!0=0$) and acts only on angular
kinetic energy.
\end{theorem}

\begin{proof}
The condensate energy functional for the gradient-knot wavefunction is
\[
  E_{\rm cond}[\psi]
  \;=\; \int \Seff(\theta_{\rm local},\rho)\,|\nabla\psi|^2\,d^3x.
\]
Decomposing via~\eqref{P13:eq:gradient_decomp} and inserting
$\Seff(0,\rho)=0$ (radial) and $\Seff(\pi/2,\rho)=\epsr$ (angular):
\[
  E_{\rm cond}[\psi]
  \;=\; \int \frac{\epsr}{r^2}\,|L\psi|^2\,d^3x
  \;=\; \left\langle\psi\left|
    -\frac{\hbar^2}{2m_e}\cdot\frac{\epsr}{r^2}\,L^2
    \right|\psi\right\rangle,
\]
where the sign is chosen so that $\Hcond$ contributes positively
to the energy (the condensate resists angular propagation).
Adding to the standard Coulomb Hamiltonian gives~\eqref{P13:eq:H_eff}.
\qed
\end{proof}

\begin{remark}[Zero radial correction]
\label{P13:rem:zero_radial}
The exact vanishing of the condensate correction for radial motion
is a structural consequence of the Hopf map Jacobian:
$\Seff(\theta)\propto\sin^4\!\theta$, and $\sin\!0=0$ exactly.
The condensate does not resist radial propagation at all.
This is not an approximation --- it is built into the geometry
of the Hopf fibration.
\end{remark}

%══════════════════════════════════════════════════════════════════════════════
\section{Chameleon Screening and the Origin of Spherical Harmonics}
\label{P13:sec:chameleon}
%══════════════════════════════════════════════════════════════════════════════

\subsection{The screened limit inside atoms}

Inside the hydrogen atom, $\bst\rhoa(r)\gg1$ throughout the
electron's orbit (equation~\eqref{P13:eq:rho_atom}).
The condensate suppression function becomes
\begin{equation}
  \epsr
  \;\approx\; \frac{1}{\vp^6\,\bst\rho_{\rm atom}(r)}
  \;\approx\; \frac{2\Lcond^4}{\vp^6\,\aem^6 m_e^4}
        \left(\frac{r}{a_0}\right)^4,
  \label{P13:eq:eps_chameleon}
\end{equation}
which is proportional to $r^4$ and vanishingly small at atomic scales.
At the Bohr radius:
\begin{equation}
  \boxed{\eps(a_0)
  \;=\; \frac{1}{\vp^6(1+\bst\rhoa(a_0))}
  \;\approx\; 2.2\times10^{-27}.}
  \label{P13:eq:eps_a0_value}
\end{equation}

\begin{theorem}[Chameleon screening: spherical harmonics from first principles]
\label{P13:thm:chameleon_screening}
In the chameleon limit $\bst\rhoa\gg1$, the condensate correction
$\Hcond = -(\hbar^2/2m_e)(\epsr/r^2)L^2$ vanishes uniformly
in all directions:
\begin{equation}
  \epsr \;\to\; 0
  \qquad \text{as } \bst\rhoa \to \infty.
  \label{P13:eq:chameleon_limit}
\end{equation}
The effective Hamiltonian reduces to the standard isotropic
Coulomb Hamiltonian:
\begin{equation}
  \Heff \;\to\; -\frac{\hbar^2}{2m_e}\nabla^2 - \frac{e^2}{r}
  \qquad (\bst\rhoa \gg 1).
  \label{P13:eq:H_standard}
\end{equation}
The eigenfunctions of~\eqref{P13:eq:H_standard} are
$\psi_{n\ell m}(r,\theta,\varphi) = R_{n\ell}(r)\,Y_\ell^m(\theta,\varphi)$,
where $Y_\ell^m$ are the standard spherical harmonics.
Spherical harmonics are therefore the correct orbital shapes
within the condensate framework, to all measurable precision.
\end{theorem}

\begin{proof}
The statement $\epsr\to0$ as $\bst\rhoa\to\infty$ follows
from~\eqref{P13:eq:epsilon_def}: as $1+\bst\rho\to\infty$,
$1/[\vp^6(1+\bst\rho)]\to0$.
The convergence is uniform in $\theta$ (note that $\epsr$
does not depend on $\theta$ --- the $\sin^4\!\theta$ anisotropy
is already extracted in~\eqref{P13:eq:H_cond}).
As $\Hcond\to0$, $\Heff\to H_{\rm Coulomb}$, whose eigenfunctions
are the standard spherical harmonic products.
\qed
\end{proof}

\begin{remark}[The same mechanism as the cosmological constant]
\label{P13:rem:same_mechanism}
The cosmological constant hierarchy is resolved in Paper~VII~\cite{Paper7}
by the chameleon mechanism: the condensate vacuum energy is present in full
but screened from gravitational coupling at cosmological density by
the denominator $1/(1+\bst\rho_\infty)$.
Theorem~\ref{P13:thm:chameleon_screening} uses the same denominator
to derive spherical harmonics inside atoms.
The two results --- the cosmological constant and orbital shapes ---
share a single algebraic origin.
\end{remark}

\begin{remark}[Quantitative screening ratio]
\label{P13:rem:screening_ratio}
The ratio of condensate density inside the atom to the cosmic attractor:
\[
  \frac{\bst\rhoa(a_0)}{\bst\rhoinf}
  \;=\; \frac{\aem^6 m_e^4}{2\Lcond^4\,\vp}
  \;\approx\; \frac{2.57\times10^{25}}{1.618}
  \;\approx\; 1.6\times10^{25}.
\]
The condensate is $1.6\times10^{25}$ times denser inside the atom
than in the cosmic vacuum.
The suppression of the condensate correction inside atoms relative
to the cosmic background is the same factor:
$\eps(a_0)/\eps_\infty \approx 10^{-27}/0.021 \approx 10^{-25}$.
\end{remark}

\subsection{Extension to multi-electron atoms}
\label{P13:sec:multi_electron}

Theorem~\ref{P13:thm:chameleon_screening} was proved for the hydrogen atom.
We now show that chameleon screening, spherical harmonics, and the bound
on condensate corrections all extend to every atom in the periodic table.

In a $Z$-electron atom, the Coulomb field at radius $r$ is generated by
the nucleus (charge $Ze$) and the electron charge distribution.
The effective nuclear charge seen at $r$ is $Z_{\rm eff}(r)$, varying
from $Z_{\rm eff}\to Z$ near the nucleus to $Z_{\rm eff}\to 1$ far from it
(where one proton is unscreened by the remaining $Z-1$ electrons).
The electromagnetic energy density is
\begin{equation}
  u_{\rm EM}^{(Z)}(r)
  \;=\; \frac{[Z_{\rm eff}(r)\,e]^2}{2r^4}
  \;=\; Z_{\rm eff}(r)^2\,u_{\rm EM}^{(1)}(r),
  \label{P13:eq:uEM_Z}
\end{equation}
and the condensate density is driven to
\begin{equation}
  \bst\rho_{\rm atom}^{(Z)}(r)
  \;\approx\; Z_{\rm eff}(r)^2\,
  \frac{\aem^6 m_e^4}{2\Lcond^4}\left(\frac{a_0}{r}\right)^4.
  \label{P13:eq:rho_Z}
\end{equation}

\begin{proposition}[Universal chameleon screening in multi-electron atoms]
\label{P13:prop:multi_electron_screening}
For any atom with $Z\ge1$ electrons, the chameleon screening condition
$\bst\rho_{\rm atom}^{(Z)}(r)\gg1$ holds for all $r$ within the
electronic density.
Consequently:
\begin{enumerate}[noitemsep,label=(\roman*)]
\item The orbital shapes are exactly the standard spherical harmonics
  $Y_\ell^m(\theta,\varphi)$ for every electron in every atom.
\item The effective Hamiltonian for the $i$-th electron reduces to
  the standard Hartree--Fock operator:
  \begin{equation}
    \Heff^{(i)} \;\to\; -\frac{\hbar^2}{2m_e}\nabla^2
    - \frac{Z_{\rm eff}(r)\,e^2}{r}
    + V_{\rm HF}^{(i)},
    \label{P13:eq:H_HF_limit}
  \end{equation}
  where $V_{\rm HF}^{(i)}$ includes the Hartree and exchange potentials.
\end{enumerate}
\end{proposition}

\begin{proof}
Equation~\eqref{P13:eq:rho_Z} gives
\begin{equation}
  \bst\rho_{\rm atom}^{(Z)}(r)
  \;\ge\; Z_{\rm eff}(r)^2\,\bst\rhoa(r)
  \;\ge\; \bst\rhoa(r)
  \;\gg\; 1,
  \label{P13:eq:rho_Z_bound}
\end{equation}
using $Z_{\rm eff}\ge1$ and the hydrogen result $\bst\rhoa(r)\gg1$
(Remark~\ref{P13:rem:radial_dep}: even at $r=5a_0$, $\bst\rhoa\approx4\times10^{22}$).
The bound $\bst\rho^{(Z)}\ge\bst\rhoa$ holds at every $r$ because the
multi-electron atom always has at least the unscreened nuclear charge
contributing to $u_{\rm EM}$.
By Theorem~\ref{P13:thm:chameleon_screening}, $\Heff^{(i)}\to H_{\rm Coulomb}^{(i)}$
whenever $\bst\rho\gg1$, giving~\eqref{P13:eq:H_HF_limit}.
\qed
\end{proof}

\begin{proposition}[Condensate energy corrections in multi-electron atoms]
\label{P13:prop:multi_electron_energy}
In the chameleon limit~\eqref{P13:eq:rho_Z_bound}, the condensate
correction to the energy of orbital $(n,\ell)$ in a $Z$-electron atom is
\begin{equation}
  \Delta E_{n\ell}^{(Z)}
  \;=\; -\frac{\hbar^2\,\ell(\ell+1)}{2m_e}
        \left\langle\frac{\eps^{(Z)}(r)}{r^2}\right\rangle_{n\ell},
  \label{P13:eq:delta_E_Z}
\end{equation}
where $\eps^{(Z)}(r)=1/[\vp^6(1+\bst\rho_{\rm atom}^{(Z)}(r))]$.
Two regimes apply:
\begin{enumerate}[noitemsep,label=(\roman*)]
\item \textbf{Inner electrons} ($Z_{\rm eff}\approx Z$, $r\sim a_0/Z$).
  In the hydrogenic approximation:
  \begin{equation}
    \Delta E_{n\ell}^{(Z,\,\rm inner)}
    \;\approx\; Z^{-4}\,\Delta E_{n\ell}^{(Z=1)},
    \label{P13:eq:delta_E_inner}
  \end{equation}
  since $\eps^{(Z)}\approx\eps^{(1)}/Z^2$ and
  $\langle r^2\rangle_{n\ell}^{(Z)}\approx\langle r^2\rangle_{n\ell}^{(1)}/Z^2$.
  The correction is \emph{more suppressed} for heavier atoms.

\item \textbf{Outer (valence) electrons} ($Z_{\rm eff}\approx 1$, $r\sim n^2 a_0$).
  Since $Z_{\rm eff}\to1$ in the outer shells,
  $\bst\rho_{\rm atom}^{(Z)}$ approaches the hydrogen value, giving:
  \begin{equation}
    \bigl|\Delta E_{n\ell}^{(Z,\,\rm outer)}\bigr|
    \;\lesssim\; \bigl|\Delta E_{n\ell}^{(Z=1)}\bigr|.
    \label{P13:eq:delta_E_outer}
  \end{equation}
\end{enumerate}
In both cases, the universal upper bound is
\begin{equation}
  \bigl|\Delta E_{n\ell}^{(Z)}\bigr|
  \;\lesssim\; \bigl|\Delta E_{n\ell}^{(Z=1)}\bigr|
  \;\sim\; 10^{-24}\,\mathrm{eV}\times\ell(\ell+1),
  \label{P13:eq:delta_E_bound}
\end{equation}
suppressed $\sim10^{20}$ below the fine structure for all atoms.
\end{proposition}

\begin{proof}
For inner electrons, the scaling $\eps^{(Z)}=\eps^{(1)}/Z^2$
follows from~\eqref{P13:eq:rho_Z} with $Z_{\rm eff}=Z$.
The hydrogenic expectation value $\langle r^2\rangle^{(Z)}=\langle r^2\rangle^{(1)}/Z^2$
(since the orbital scale is $a_0/Z$).
Combining: $\langle\eps^{(Z)}/r^2\rangle^{(Z)}=Z^{-4}\langle\eps^{(1)}/r^2\rangle^{(1)}$,
giving~\eqref{P13:eq:delta_E_inner}.
For outer electrons, $Z_{\rm eff}\approx1$ restores the hydrogen formula,
giving~\eqref{P13:eq:delta_E_outer}.
The bound~\eqref{P13:eq:delta_E_bound} follows since
$|\Delta E_{n\ell}^{(Z=1)}|\approx1.77\times10^{-24}\,\mathrm{eV}$
for the $2p$ state (Proposition~\ref{P13:prop:energy_correction}).
\qed
\end{proof}

\begin{corollary}[The Hartree--Fock equations are exact]
\label{P13:cor:HF_exact}
Proposition~\ref{P13:prop:multi_electron_screening} establishes that the
condensate correction to the Hartree--Fock Hamiltonian is
$O(\eps^{(Z)})\lesssim O(10^{-27})$, negligible at any measured precision.
Standard quantum chemistry methods --- Hartree--Fock, DFT, coupled cluster,
configuration interaction --- are therefore all derivable within the
condensate framework, with no condensate correction needed.
In particular, the Aufbau filling order and the structure of the
periodic table (addressed via the $\vp$-spiral in Paper~XI~\cite{Paper11})
are unaffected by the condensate.
\end{corollary}

\begin{theorem}[Atomic physics from the condensate]
\label{P13:thm:atomic_foundation}
This paper, taken together, derives from the condensate framework:
\begin{enumerate}[noitemsep]
\item The modified Hamiltonian with a single condensate correction
  (Theorem~\ref{P13:thm:modified_H}).
\item Spherical harmonics as exact orbital shapes for all atoms and all $Z$
  (Theorem~\ref{P13:thm:chameleon_screening} and
  Proposition~\ref{P13:prop:multi_electron_screening}).
\item Condensate energy corrections bounded by $\sim10^{-24}\,\mathrm{eV}$,
  undetectable in any atom (Propositions~\ref{P13:prop:energy_correction}
  and~\ref{P13:prop:multi_electron_energy}).
\item The Schr\"odinger equation from condensate quantisation,
  exact to $O(\eps_{\rm cond})\approx10^{-27}$
  (Theorem~\ref{P13:thm:schrodinger}, Corollary~\ref{P13:cor:improved_BO}).
\item Pauli exclusion and spin-$\tfrac{1}{2}$ from $Q=2$ topology
  (Theorems~\ref{P13:thm:spin_half}--\ref{P13:cor:pauli_exclusion}).
\item The spin assignments $J=Q/4$ for all Standard Model particles
  (Theorem~\ref{P13:thm:spin_tower}).
\end{enumerate}
Standard atomic and molecular quantum mechanics --- wavefunctions,
selection rules, the periodic table, all of quantum chemistry ---
follows from the condensate framework without additional postulates.
\end{theorem}

%══════════════════════════════════════════════════════════════════════════════
\section{Condensate Corrections to Hydrogen Energy Levels}
\label{P13:sec:energy_correction}
%══════════════════════════════════════════════════════════════════════════════

\begin{proposition}[Condensate energy correction to hydrogen levels]
\label{P13:prop:energy_correction}
In first-order perturbation theory with $H' = \Hcond$,
the condensate correction to the hydrogen energy level $|n,\ell\rangle$ is
\begin{equation}
  \Delta E_{n\ell}
  \;=\; \left\langle n,\ell\left|H'\right|n,\ell\right\rangle
  \;=\; -\frac{\hbar^2\ell(\ell+1)}{2m_e}\,
        \left\langle\frac{\epsr}{r^2}\right\rangle_{n\ell}.
  \label{P13:eq:delta_E_general}
\end{equation}
In the chameleon limit $\bst\rhoa\gg1$, inserting~\eqref{P13:eq:eps_chameleon}:
\begin{equation}
  \Delta E_{n\ell}
  \;\approx\; -\frac{\ell(\ell+1)\,\Lcond^4}{\vp^6\,\aem^4\,m_e^3}
              \cdot\frac{\langle r^2\rangle_{n\ell}}{a_0^4},
  \label{P13:eq:delta_E_chameleon}
\end{equation}
where $\langle r^2\rangle_{n\ell} = (n^2/2)(5n^2+1-3\ell(\ell+1))\,a_0^2$
is the exact hydrogen expectation value.
For the $2p$ state ($n=2$, $\ell=1$):
\begin{equation}
  \boxed{\Delta E_{2p}
  \;\approx\; -\frac{2\,\Lcond^4\cdot30}{\vp^6\,\aem^4\,m_e^3}
  \;\approx\; -1.77\times10^{-24}\,\mathrm{eV}.}
  \label{P13:eq:delta_E_2p}
\end{equation}
This is $\sim1.95\times10^{-20}$ times the fine-structure
splitting of the $2p$ level
($\Delta E_{\rm fs}\sim\aem^2\ERy/n^3\approx9.1\times10^{-5}\,\mathrm{eV}$).
\end{proposition}

\begin{proof}
Equation~\eqref{P13:eq:delta_E_general} is immediate from first-order
perturbation theory applied to $H'=\Hcond$,
using $L^2|n,\ell,m\rangle = \ell(\ell+1)\hbar^2|n,\ell,m\rangle$.
For~\eqref{P13:eq:delta_E_chameleon}: in the chameleon limit,
$\epsr \approx 2\Lcond^4 r^4/[\vp^6\aem^6 m_e^4 a_0^4]$
(from~\eqref{P13:eq:eps_chameleon}), so
$\langle\epsr/r^2\rangle_{n\ell}
= 2\Lcond^4/[\vp^6\aem^6 m_e^4 a_0^4]\cdot\langle r^2\rangle_{n\ell}$.
Inserting the standard hydrogen result
$\langle r^2\rangle_{2p} = 30a_0^2$ and simplifying
using $\hbar^2/(2m_e a_0^2) = \aem^2 m_e/2 = \ERy$
gives~\eqref{P13:eq:delta_E_2p}.
\qed
\end{proof}

\begin{remark}[Structural form of the correction]
\label{P13:rem:structure}
The condensate correction~\eqref{P13:eq:delta_E_general}:
\begin{enumerate}[noitemsep]
\item Is zero for $\ell=0$ states (s-orbitals), exactly:
  $L^2|n,0\rangle = 0$.
\item Breaks the $\ell$-degeneracy of hydrogen energy levels
  (the SO(4) accidental degeneracy): states with the same $n$
  but different $\ell$ receive different corrections.
  This is a condensate analogue of the Lamb shift, at $\sim10^{-20}$
  smaller magnitude.
\item Scales as $\Lcond^4/m_e^3 \propto \TCMB^4/m_e^3$:
  it would grow as the CMB temperature increases, meaning the
  condensate correction was larger in the early universe.
\end{enumerate}
\end{remark}

%══════════════════════════════════════════════════════════════════════════════
\section{Free-Electron Mass Anisotropy in the Cosmic Background}
\label{P13:sec:anisotropy}
%══════════════════════════════════════════════════════════════════════════════

Outside atoms, in the cosmic background at the condensate attractor
density $\rhoinf = \vp/\bst$, the condensate is NOT chameleon-screened.
The condensate correction is:
\begin{equation}
  \boxed{\eps(\rhoinf)
  \;=\; \frac{1}{\vp^6(1+\bst\rhoinf)}
  \;=\; \frac{1}{\vp^6\cdot\vp^2}
  \;=\; \frac{1}{\vp^8},}
  \label{P13:eq:eps_attractor}
\end{equation}
where we used $\bst\rhoinf=\vp$ (attractor definition)
and $1+\vp=\vp^2$ (exact pentagon identity).
This is non-negligible.

\begin{proposition}[Free-electron mass anisotropy]
\label{P13:prop:mass_anisotropy}
For a free electron propagating in the cosmic condensate background
at the attractor density $\rhoinf$, the effective mass depends on
the direction of propagation relative to the local condensate gradient:
\begin{align}
  m_{\rm radial}
  &= m_e
  &\text{(exactly, for propagation parallel to } \nabla\rho\text{)},
  \label{P13:eq:m_radial}
  \\
  m_{\rm tangential}
  &= m_e\left(1 + \frac{1}{\vp^8}\right)
   = m_e\,\frac{\vp^8+1}{\vp^8}
  &\text{(perpendicular to } \nabla\rho\text{)}.
  \label{P13:eq:m_tangential}
\end{align}
The fractional mass anisotropy is
\begin{equation}
  \boxed{\frac{\Delta m}{m_e}
  \;=\; \frac{m_{\rm tangential} - m_{\rm radial}}{m_e}
  \;=\; \frac{1}{\vp^8}
  \;\approx\; 2.129\%.}
  \label{P13:eq:mass_anisotropy}
\end{equation}
The radial-to-tangential effective mass ratio is
\begin{equation}
  \frac{m_{\rm radial}}{m_{\rm tangential}}
  \;=\; \frac{\vp^8}{\vp^8+1}
  \;\approx\; 0.97916,
  \label{P13:eq:mass_ratio}
\end{equation}
exact in terms of $\vp$ with no free parameters.
\end{proposition}

\begin{proof}
Equations~\eqref{P13:eq:m_radial}--\eqref{P13:eq:m_tangential} follow from
Theorem~\ref{P13:thm:modified_H} applied to a free electron
(no Coulomb potential) with the cosmic background density $\rhoinf$.
The condensate correction to the kinetic energy for motion in
direction $\hat{e}$ at angle $\theta$ to $\nabla\rho$ is
$\eps(\rhoinf)\sin^4\!\theta\,p^2/(2m_e)$, giving effective mass
$m_{\rm eff}(\theta) = m_e/(1+\eps(\rhoinf)\sin^4\!\theta)$.
Evaluating at $\theta=0$ and $\theta=\pi/2$ with~\eqref{P13:eq:eps_attractor}
gives~\eqref{P13:eq:m_radial} and~\eqref{P13:eq:m_tangential}
(at linear order in $\eps\approx0.021\ll1$, the inverse is
$1/(1+\eps\sin^4\!\theta)\approx1-\eps\sin^4\!\theta$,
but the exact form is as stated).
The numerical value $1/\vp^8 = 0.021286\ldots$ follows
from $\vp^8 = 46.9787\ldots$\,.
\qed
\end{proof}

\begin{remark}[Exact form and approximations]
\label{P13:rem:exact_vs_approx}
The result~\eqref{P13:eq:mass_anisotropy} uses two exact golden-ratio
identities: $\bst\rhoinf=\vp$ (the condensate self-consistency
condition, Paper~I~\cite{Paper1}) and $1+\vp=\vp^2$
(the pentagon identity).
The mass ratio $\vp^8/(\vp^8+1)$ involves no free parameters.
The $2.129\%$ anisotropy follows from these alone.
\end{remark}

\begin{remark}[Direction of the preferred axis and experimental access]
\label{P13:rem:experimental_access}
The condensate gradient $\nabla\rho$ near a gravitational source
(Earth, Sun, Galaxy) is radially directed.
In empty space far from any mass, the large-scale condensate
gradient is determined by the local matter distribution and the
CMB dipole.
A free electron (or any lepton) propagating radially versus
tangentially relative to the galactic centre, or relative to
the CMB dipole direction, should show a $2.129\%$ difference
in effective mass.
The natural experimental probe is a precision Penning trap:
the cyclotron frequency $\omega_c = eB/m_{\rm eff}$ would depend
on the orientation of the trap relative to the local condensate
gradient at the $2.129\%/\vp^8$ level.

Current Penning-trap measurements of the electron $g$-factor
(Hanneke et al.~\cite{Hanneke2008}) have fractional precision
$\sim10^{-13}$, which could in principle detect a $2\%$ effect
if the anisotropy is correctly modelled.
A dedicated anisotropy search --- varying trap orientation
over a year to scan the local condensate gradient direction ---
would provide a clean test.
\end{remark}

\begin{remark}[Relation to the $\alpha$--$m_e$ tracking relation]
\label{P13:rem:tracking_relation}
Paper~IV~\cite{Paper4} establishes the tracking relation
$\delta m_e/m_e = -\delta\aem/(\pi\aem) \approx -43.6\,\delta\aem/\aem$:
as $\aem$ is refined, $m_e$ tracks.
The mass anisotropy~\eqref{P13:eq:mass_anisotropy} is a distinct
effect: the anisotropy is in the EFFECTIVE MASS during propagation,
not in the rest mass $m_e$.
The rest mass is set by the tower level and WZW T-phase (Papers~I, IV)
and is isotropic.
The propagation effective mass acquires the $1/\vp^8$ anisotropy
from the condensate medium.
\end{remark}

%══════════════════════════════════════════════════════════════════════════════
\section{The Schr\"odinger Equation from the Condensate Quantum Theory}
\label{P13:sec:schrodinger}
%══════════════════════════════════════════════════════════════════════════════

Paper~X~\cite{Paper10} carries out the canonical quantisation of the
condensate field, producing the Hilbert space
$\mathcal{H}=L^2(\mathcal{C}_Q^{\rm red},d\mu)$, the Born rule,
and the golden-ratio spectrum.
Sections~\ref{P13:sec:modified_H}--\ref{P13:sec:anisotropy} of this paper
used the Schr\"odinger equation as given.
We now derive it from Paper~X's quantum theory,
using no additional postulates beyond those already in the condensate framework.

\subsection{The one-particle sector and the translation fibration}
\label{P13:sec:translation_fibration}

The condensate field configuration $\rho(x)$ describing a single
gradient-knot (electron) localised near position $\mathbf{R}\in\mathbb{R}^3$
can be written as $\rho(x) = \rho_{\rm knot}(x-\mathbf{R})$,
where $\rho_{\rm knot}$ is the profile of the knot centred at the origin.

The group $\mathbb{R}^3$ acts on $\mathcal{C}_Q^{\rm red}$ by translations:
\begin{equation}
  T_{\mathbf{a}}[\rho](x) = \rho(x-\mathbf{a}),
  \label{P13:eq:translation_action}
\end{equation}
and this action is free on the one-knot sector.
The reduced configuration space therefore fibres over position space:
\begin{equation}
  \mathcal{C}_Q^{\rm red} \;=\; \mathbb{R}^3 \times \mathcal{C}_{\rm internal},
  \label{P13:eq:fibration}
\end{equation}
where $\mathcal{C}_{\rm internal} = \mathcal{C}_Q^{\rm red}/\mathbb{R}^3$
is the space of internal shapes (profile functions
$f(r,z)$ modulo translation).

The Liouville measure $d\mu$ on $\mathcal{C}_Q^{\rm red}$
(Paper~X~\cite{Paper10}, Corollary~8.4) factorises:
$d\mu = d^3R\otimes d\mu_{\rm internal}$.
The Hilbert space accordingly decomposes:
\begin{equation}
  \mathcal{H}
  \;=\; L^2\!\left(\mathcal{C}_Q^{\rm red},d\mu\right)
  \;\cong\;
  L^2\!\left(\mathbb{R}^3,d^3R\right)\otimes\mathcal{H}_{\rm internal},
  \label{P13:eq:Hilbert_factorisation}
\end{equation}
where $\mathcal{H}_{\rm internal} = L^2(\mathcal{C}_{\rm internal},d\mu_{\rm internal})$.

\begin{definition}[One-particle states]
\label{P13:def:one_particle}
A one-particle state is a state of the form
\begin{equation}
  |\Psi_e\rangle \;=\; \psi(\mathbf{R})\otimes|0_{\rm internal}\rangle,
  \label{P13:eq:one_particle_state}
\end{equation}
where $\psi\in L^2(\mathbb{R}^3)$ is the centre-of-mass wavefunction
and $|0_{\rm internal}\rangle\in\mathcal{H}_{\rm internal}$
is the ground state of the condensate knot's internal structure.
\end{definition}

\subsection{The momentum operator from the Noether current}
\label{P13:sec:momentum_operator}

The parent $k$-essence Lagrangian (Paper~I~\cite{Paper1})
\begin{equation}
  \mathcal{L} \;=\; \frac{|\nabla\rho|^2}{\vp^6(1+\bst|\nabla\rho|^2)}
  \label{P13:eq:parent_L}
\end{equation}
is invariant under spatial translations~\eqref{P13:eq:translation_action}.
By Noether's theorem, the conserved momentum density is
\begin{equation}
  \mathcal{P}_i(x)
  \;=\; \pi(x)\,\partial_i\rho(x),
  \qquad
  \pi(x) \;=\; \frac{\partial\mathcal{L}}{\partial(\partial_0\rho)},
  \label{P13:eq:noether_momentum}
\end{equation}
and the total conserved spatial momentum is
\begin{equation}
  P_i \;=\; \int \pi(x)\,\partial_i\rho(x)\,d^3x.
  \label{P13:eq:total_momentum}
\end{equation}
This is the momentum of Paper~X~\cite{Paper10} (Lemma~lem:momentum),
proved to generate the symplectic structure on $\mathcal{C}_Q^{\rm red}$.

\begin{lemma}[Momentum operator on the one-particle sector]
\label{P13:lem:momentum_operator}
On the one-particle sector~\eqref{P13:eq:one_particle_state},
the operator corresponding to the Noether momentum~\eqref{P13:eq:total_momentum}
acts as
\begin{equation}
  \hat{P}_i\,\psi(\mathbf{R})\otimes|0_{\rm internal}\rangle
  \;=\; \bigl(-i\hbar\,\partial/\partial R_i\,\psi(\mathbf{R})\bigr)
        \otimes|0_{\rm internal}\rangle.
  \label{P13:eq:momentum_on_sector}
\end{equation}
\end{lemma}

\begin{proof}
On the one-particle sector, $T_{\mathbf{a}}$ acts as
$\psi(\mathbf{R})\mapsto\psi(\mathbf{R}+\mathbf{a})$.
By Stone's theorem, the generator of this unitary one-parameter group
is $\hat{P}_i = -i\hbar\,\partial/\partial R_i$
on $L^2(\mathbb{R}^3)$.
The internal state $|0_{\rm internal}\rangle$ is translation-invariant
(it is the ground state of the shape space), so $P_i$
acts only on the $\psi$ factor.
\qed
\end{proof}

\subsection{The effective mass and the Born--Oppenheimer separation}
\label{P13:sec:born_oppenheimer}

The full condensate Hamiltonian on $\mathcal{H}$ decomposes as
\begin{equation}
  \hat{H}_{\rm cond}
  \;=\; \hat{H}_{\rm cm}(\mathbf{R})
  \;+\; \hat{H}_{\rm internal}
  \;+\; \hat{V}(\mathbf{R}),
  \label{P13:eq:H_decomp}
\end{equation}
where $\hat{H}_{\rm cm}$ is the centre-of-mass kinetic energy,
$\hat{H}_{\rm internal}$ acts only on $\mathcal{H}_{\rm internal}$,
and $\hat{V}$ is any external potential (e.g.\ the Coulomb potential
of the proton).

The key scale separation is between:
\begin{itemize}[noitemsep]
  \item \emph{Electron kinetic energy at the Bohr radius:}
    $\ERy = \aem^2 m_e/2 \approx 13.6\,\mathrm{eV}$.
  \item \emph{First internal excitation of the electron gradient-knot:}
    The knot is at tower level $n=20$; its first excitation is at level $n=19$,
    with energy gap
    \begin{equation}
      \boxed{\Delta_{\rm internal}
      \;=\; m_e(\vp^{-38}-\vp^{-40})\big/\vp^{-40}
      \;=\; m_e(\vp^2-1)
      \;=\; \vp\, m_e
      \;\approx\; 827\,\mathrm{keV},}
      \label{P13:eq:internal_gap}
    \end{equation}
    using $\vp^2-1=\vp$ (pentagon identity).
\end{itemize}

The Born--Oppenheimer adiabatic parameter is therefore
\begin{equation}
  \boxed{\eps_{\rm BO}
  \;=\; \frac{\ERy}{\vp\, m_e}
  \;=\; \frac{\aem^2}{2\vp}
  \;\approx\; 1.6\times10^{-5}.}
  \label{P13:eq:eps_BO}
\end{equation}

\begin{theorem}[Born--Oppenheimer separation for the condensate]
\label{P13:thm:born_oppenheimer}
In the adiabatic limit $\eps_{\rm BO}\to0$, the condensate Hamiltonian~\eqref{P13:eq:H_decomp}
restricted to one-particle states~\eqref{P13:eq:one_particle_state} with
$|0_{\rm internal}\rangle$ in the ground state of $\hat{H}_{\rm internal}$ becomes
\begin{equation}
  \hat{H}_{\rm eff}
  \;=\; \frac{\hat{\mathbf{P}}^2}{2m_e}
  \;+\; \hat{V}(\mathbf{R})
  \;+\; O\!\left(\eps_{\rm BO}^2\right),
  \label{P13:eq:Heff_BO}
\end{equation}
where $m_e = \Lcond\,\vp^{20}\,e^{4\pi-3/400-\aem/\pi}$
is the condensate-derived electron mass (Papers~I,~IV~\cite{Paper1,Paper4}).
\end{theorem}

\begin{proof}
\textbf{Step 1 (Effective mass).}
For a knot translating with velocity $\mathbf{v}$,
the centre-of-mass kinetic energy is
\begin{equation}
  E_{\rm cm}
  \;=\; \int\pi(x)\,\partial_0\rho(x)\,d^3x - \mathcal{L}
  \;\approx\; \tfrac{1}{2}M_{\rm eff}\,\mathbf{v}^2,
  \label{P13:eq:E_cm}
\end{equation}
in the non-relativistic limit $|\mathbf{v}|\ll c$.
The effective mass $M_{\rm eff}$ is the condensate's
inertia to translation.
For the $Q=10$ gradient-knot at tower level $n=20$,
Paper~I~\cite{Paper1} fixes the static energy as
$E_{\rm static} = \mathcal{K}_{\rm fb}\cdot\mathcal{J}_4 = m_e c^2$
(equation~thm:yukawa of Paper~I).
Lorentz covariance of the parent
Lagrangian~\eqref{P13:eq:parent_L} then requires $M_{\rm eff}=E_{\rm static}/c^2=m_e$:
the effective mass equals the condensate rest mass.

\textbf{Step 2 (Adiabatic theorem).}
The energy gap~\eqref{P13:eq:internal_gap} is $\Delta_{\rm internal}=\vp\,m_e$.
For the electron moving at atomic energies $E_{\rm cm}\sim\ERy$,
the condition $E_{\rm cm}\ll\Delta_{\rm internal}$ is satisfied to
$\eps_{\rm BO} = \aem^2/(2\vp)\approx1.6\times10^{-5}$.
By the quantum adiabatic theorem, the internal state $|0_{\rm internal}\rangle$
is preserved during the knot's motion to $O(\eps_{\rm BO}^2)$.

\textbf{Step 3 (Effective Hamiltonian).}
With the internal state frozen, the effective Hamiltonian on
$L^2(\mathbb{R}^3)$ is
$\hat{H}_{\rm eff}
= \langle 0_{\rm internal}|\hat{H}_{\rm cm}|0_{\rm internal}\rangle+\hat{V}
= \hat{\mathbf{P}}^2/(2m_e)+\hat{V}$,
where the last step uses $M_{\rm eff}=m_e$ from Step~1
and Lemma~\ref{P13:lem:momentum_operator}.
The correction is the non-adiabatic coupling
$\langle 0_{\rm internal}|\nabla_\mathbf{R}|0_{\rm internal}\rangle=0$
(vanishes by time-reversal symmetry of the ground state),
so the leading correction is $O(\eps_{\rm BO}^2)$.
\qed
\end{proof}

\subsection{The Schr\"odinger equation}
\label{P13:sec:schrodinger_eq}

\begin{theorem}[Schr\"odinger equation from condensate quantisation]
\label{P13:thm:schrodinger}
In the one-particle sector of the condensate Hilbert space,
the time evolution of the electron's centre-of-mass wavefunction
$\psi(\mathbf{R},t)\in L^2(\mathbb{R}^3)$ satisfies
\begin{equation}
  i\hbar\,\frac{\partial\psi}{\partial t}
  \;=\;
  \left[-\frac{\hbar^2}{2m_e}\nabla^2 + V(\mathbf{R})\right]\psi
  \;+\; O\!\left(\eps_{\rm cond}\right),
  \label{P13:eq:schrodinger_derived}
\end{equation}
where $m_e$ is the condensate-derived electron mass,
$V(\mathbf{R})$ is any external potential,
and $\eps_{\rm cond}\sim(\Lcond/m_e\aem)^4\sim10^{-27}$ is the
condensate correction from Theorem~\ref{P13:thm:modified_H}.
The $O(\eps_{\rm BO}^2)$ Born--Oppenheimer remainder vanishes exactly
by the translation fibration (Theorem~\ref{P13:thm:NA_vanishing}
and Corollary~\ref{P13:cor:improved_BO}); the error bound
$O(\eps_{\rm cond})\approx O(10^{-27})$ is a factor of
$\sim10^{17}$ tighter than the formal BO bound.
\end{theorem}

\begin{proof}
Paper~X~\cite{Paper10}, Theorem~thm:quantisation provides the time-evolution
equation $i\hbar\,\partial|\Psi\rangle/\partial t = \hat{H}|\Psi\rangle$
on the full condensate Hilbert space $\mathcal{H}$.
Projecting onto the one-particle sector via Definition~\ref{P13:def:one_particle},
applying Theorem~\ref{P13:thm:born_oppenheimer} to replace $\hat{H}$
by $\hat{H}_{\rm eff} = \hat{\mathbf{P}}^2/(2m_e)+\hat{V}$,
and using $\hat{P}_i = -i\hbar\partial/\partial R_i$
(Lemma~\ref{P13:lem:momentum_operator}), gives
\begin{equation*}
  i\hbar\,\frac{\partial\psi}{\partial t}
  = \left[-\frac{\hbar^2}{2m_e}\nabla_\mathbf{R}^2 + V(\mathbf{R})\right]\psi
  + O(\eps_{\rm BO}^2).
\end{equation*}
Adding the condensate anisotropy correction $\Hcond$ from
Theorem~\ref{P13:thm:modified_H}, suppressed to $O(\eps_{\rm cond})$
by Theorem~\ref{P13:thm:chameleon_screening}, gives~\eqref{P13:eq:schrodinger_derived}.
\qed
\end{proof}

\begin{remark}[What the derivation establishes]
\label{P13:rem:what_derived}
Theorem~\ref{P13:thm:schrodinger} shows that the Schr\"odinger equation
is not an additional postulate of the condensate framework.
It is a consequence of four structural results already present:
\begin{enumerate}[noitemsep]
  \item The Hilbert space $\mathcal{H}=L^2(\mathcal{C}_Q^{\rm red},d\mu)$
    and time-evolution equation (Paper~X~\cite{Paper10}, Theorem~thm:quantisation).
  \item The translation fibration $\mathcal{C}_Q^{\rm red}=\mathbb{R}^3\times
    \mathcal{C}_{\rm internal}$ (this section, equation~\eqref{P13:eq:fibration}).
  \item The Noether momentum $\hat{P}_i=-i\hbar\partial/\partial R_i$
    on the one-particle sector (Lemma~\ref{P13:lem:momentum_operator}).
  \item The Born--Oppenheimer adiabatic separation at
    $\eps_{\rm BO}=\aem^2/(2\vp)\approx1.6\times10^{-5}$
    (Theorem~\ref{P13:thm:born_oppenheimer}).
\end{enumerate}
The electron mass $m_e$ appearing in~\eqref{P13:eq:schrodinger_derived}
is not an input to this derivation: it is the condensate mass
from Papers~I and~IV~\cite{Paper1,Paper4}, derived from $\TCMB$ alone.
\end{remark}

\begin{remark}[Adiabatic correction estimate — initial bound]
\label{P13:rem:BO_correction}
Theorem~\ref{P13:thm:born_oppenheimer} states the Born--Oppenheimer error as
$O(\eps_{\rm BO}^2)$ where $\eps_{\rm BO}=\aem^2/(2\vp)\approx1.6\times10^{-5}$,
giving a formal bound of $\eps_{\rm BO}^2\approx2.7\times10^{-10}$
on the fractional energy correction.
This bound is sharp for a generic BO problem (e.g.\ a diatomic molecule),
where the derivative coupling $d_{0k}=\langle k|\partial/\partial R|0\rangle\neq0$.
We now show that in the condensate framework this coupling vanishes
\emph{identically}, improving the bound by a factor of $\sim10^{17}$.
\end{remark}

\begin{theorem}[Exact vanishing of non-adiabatic coupling]
\label{P13:thm:NA_vanishing}
The Born--Oppenheimer derivative coupling between any two internal states
of the $Q=2$ gradient-knot vanishes identically:
\begin{equation}
  \boxed{d_{0k} \;\equiv\; \langle k_{\rm internal}|\,\partial/\partial R_i\,|0_{\rm internal}\rangle
  \;=\; 0
  \qquad \text{for all } k \text{ and all } i=1,2,3.}
  \label{P13:eq:d0k_zero}
\end{equation}
Consequently, the diagonal Born--Oppenheimer correction (DBOC) also vanishes:
\begin{equation}
  \mathrm{DBOC}
  \;\equiv\; -\frac{\hbar^2}{2m_e}\sum_{k\neq0}\frac{|d_{0k}|^2}{E_k-E_0}
  \;=\; 0
  \label{P13:eq:DBOC_zero}
\end{equation}
and the $O(\eps_{\rm BO}^2)$ term in Theorem~\ref{P13:thm:schrodinger}
is exactly zero.
\end{theorem}

\begin{proof}
In standard molecular Born--Oppenheimer theory, the derivative coupling
$d_{0k}\neq0$ because the electronic wavefunction $\Phi_0(r;R)$
depends explicitly on the nuclear coordinate $R$: as the nucleus moves,
the electrons follow, generating a non-zero gradient $\partial\Phi_0/\partial R$.

In the condensate framework, the situation is structurally different.
From the translation fibration~\eqref{P13:eq:fibration}:
\begin{equation}
  \mathcal{C}_Q^{\rm red} \;=\; \mathbb{R}^3 \times \mathcal{C}_{\rm internal},
  \label{P13:eq:fibration_recalled}
\end{equation}
the internal configuration space $\mathcal{C}_{\rm internal}=\mathcal{C}_Q^{\rm red}/\mathbb{R}^3$
is defined by modding out ALL translations.
The internal states
$|k_{\rm internal}\rangle\in\mathcal{H}_{\rm internal}=L^2(\mathcal{C}_{\rm internal},d\mu_{\rm internal})$
therefore have \emph{no $R$-dependence by construction}:
\begin{equation}
  \frac{\partial}{\partial R_i}\,|k_{\rm internal}\rangle \;=\; 0
  \qquad\text{for all } k \text{ and all } R.
  \label{P13:eq:internal_R_indep}
\end{equation}
This holds to all orders; it is not an approximation.
Equation~\eqref{P13:eq:d0k_zero} follows immediately, and
equation~\eqref{P13:eq:DBOC_zero} follows from~\eqref{P13:eq:d0k_zero}.

In physical terms: the ``electrons'' in standard BO are tied to the nuclear position
by the Coulomb potential, so they move when the nucleus moves.
The condensate internal states are in the \emph{shape space} modulo translations;
they are defined to be position-independent and cannot couple to the center-of-mass motion.
\qed
\end{proof}

\begin{corollary}[Improved error bound for the Schr\"odinger equation]
\label{P13:cor:improved_BO}
The Schr\"odinger equation~\eqref{P13:eq:schrodinger_derived} of
Theorem~\ref{P13:thm:schrodinger} holds with the improved error:
\begin{equation}
  \boxed{i\hbar\,\frac{\partial\psi}{\partial t}
  \;=\;
  \left[-\frac{\hbar^2}{2m_e}\nabla^2 + V(\mathbf{R})\right]\psi
  \;+\; O(\eps_{\rm cond}),}
  \label{P13:eq:schrodinger_improved}
\end{equation}
where $\eps_{\rm cond}\sim10^{-27}$ is the condensate anisotropy
from Theorem~\ref{P13:thm:modified_H}.
The formal $O(\eps_{\rm BO}^2)$ term vanishes exactly
(Theorem~\ref{P13:thm:NA_vanishing}), improving the error bound by
a factor of
\begin{equation}
  \boxed{\frac{\eps_{\rm BO}^2}{\eps_{\rm cond}}
  \;=\; \frac{[\aem^2/(2\vp)]^2}{\eps_{\rm cond}}
  \;\approx\; \frac{2.7\times10^{-10}}{2.2\times10^{-27}}
  \;\approx\; 1.2\times10^{17}.}
  \label{P13:eq:improvement_factor}
\end{equation}
The Schr\"odinger equation is therefore exact to one part in $10^{27}$,
not merely one part in $10^{10}$.
\end{corollary}

\begin{proof}
Theorem~\ref{P13:thm:NA_vanishing} shows $d_{0k}=0$ and DBOC$=0$,
so the $O(\eps_{\rm BO}^2)$ remainder in Theorem~\ref{P13:thm:born_oppenheimer}
is absent.
The only remaining correction is the condensate anisotropy
$\Hcond=O(\eps_{\rm cond})$ from Theorem~\ref{P13:thm:modified_H},
which is already included in~\eqref{P13:eq:schrodinger_derived}.
\qed
\end{proof}

\begin{remark}[Why the condensate BO is exact and molecular BO is not]
\label{P13:rem:BO_comparison}
The contrast between the condensate and molecular cases is sharp:
\begin{enumerate}[noitemsep]
\item \emph{Molecule} (e.g.\ H$_2$):
  Electronic wavefunction $\Phi_0(r;R)$ depends on nuclear coordinate $R$.
  Derivative coupling $d_{0k}\neq0$ ($\sim\sqrt{m_e/M}$ per matrix element).
  Non-adiabatic correction $\sim\eps_{\rm BO}\sim m_e/M$.
\item \emph{Condensate gradient-knot} (this paper):
  Internal states in $\mathcal{H}_{\rm internal}$ are $R$-independent by the
  translation fibration~\eqref{P13:eq:fibration}.
  Derivative coupling $d_{0k}=0$ \emph{exactly} for all $k$.
  Non-adiabatic correction $=0$; residual error $=O(\eps_{\rm cond})$.
\end{enumerate}
The condensate BO is \emph{exact} in a sense that ordinary molecular BO is not:
the internal and external degrees of freedom are structurally decoupled
by the configuration-space fibration, not merely approximately so by a large mass ratio.
\end{remark}

\begin{remark}[Condensate phonons and the electron]
\label{P13:rem:phonons}
The condensate background has phonon excitations at
$\Delta_{\rm phonon}\sim\vp^2\Lcond\approx3\times10^{-4}\,\mathrm{eV}$,
which are \emph{below} the electron kinetic energy $\ERy\approx13.6\,\mathrm{eV}$.
The electron can in principle excite condensate phonons as it orbits.
However, the coupling of the electron gradient-knot to condensate
phonons is suppressed by the same chameleon factor
$\eps_{\rm cond}\sim10^{-27}$ (Section~\ref{P13:sec:chameleon}):
the atom's EM field density overwhelms the condensate background density,
screening all condensate interactions.
Net phonon emission rate: $\Gamma_{\rm phonon}\sim\eps_{\rm cond}^2\times
\text{(phase space)}\ll10^{-50}\,\mathrm{s}^{-1}$.
The electron's quantum mechanics is entirely standard at atomic energies.
\end{remark}

%══════════════════════════════════════════════════════════════════════════════
\section{Fermionic Statistics and the Pauli Exclusion Principle}
\label{P13:sec:fermi}
%══════════════════════════════════════════════════════════════════════════════

The preceding sections referred to the electron as a fermionic object
without deriving this property from the condensate framework.
We show here that Fermi statistics and the Pauli exclusion principle
are not additional postulates: they follow directly from the icosahedral
symmetry and the Hopf charge $Q=2$ already established in Paper~I~\cite{Paper1}.

\subsection{The binary icosahedral group as a double cover}
\label{P13:sec:2I_double_cover}

Paper~I establishes that the condensate has icosahedral symmetry,
with symmetry group the binary icosahedral group $2I$ (also written $2I\cong\tilde{I}$).
The group $2I$ is a subgroup of $SU(2)$ and a double cover of the
ordinary icosahedral group $I\subset SO(3)$:
\begin{equation}
  1 \;\longrightarrow\; \mathbb{Z}_2
  \;\longrightarrow\; 2I
  \;\stackrel{\pi}{\longrightarrow}\; I
  \;\longrightarrow\; 1,
  \label{P13:eq:2I_exact_seq}
\end{equation}
where the kernel of $\pi$ is the centre $Z(2I)=\mathbb{Z}_2=\{e,z\}$.
The central element $z$ satisfies $z^2=e$, $z\neq e$, and represents the
$2\pi$ spatial rotation: it is the image of $-\mathbf{1}\in SU(2)$
under the standard embedding $2I\subset SU(2)$.

The irreducible representations of $2I$ fall into two classes:
\begin{itemize}[noitemsep]
  \item \emph{Integer-spin} (bosonic): representations in which
    $z$ acts as $+\mathbf{1}$, i.e.\ they descend to representations
    of $I$.
    Dimensions: 1, 3, 4, 5 (from the 5 conjugacy classes of $I$).
  \item \emph{Half-integer-spin} (fermionic): representations in which
    $z$ acts as $-\mathbf{1}$, peculiar to the double cover.
    Dimensions: 2, 2, 4, 6.
    These do not descend to representations of $I$.
\end{itemize}
The two-dimensional half-integer representation (the fundamental spinor
of $2I\subset SU(2)$) is the $j=\tfrac{1}{2}$ representation of $SU(2)$
restricted to $2I$.
Under the central element $z$ (i.e.\ a $2\pi$ rotation):
\begin{equation}
  \hat{T}_z \,\psi = z\cdot\psi = (-1)\cdot\psi = -\psi.
  \label{P13:eq:z_action}
\end{equation}

\begin{theorem}[Spin-$\tfrac{1}{2}$ of the Q=2 gradient-knot]
\label{P13:thm:spin_half}
The electron gradient-knot (Hopf charge $Q=2$, Paper~I) transforms
in the fundamental spinor representation of $2I$ and carries spin $J=\tfrac{1}{2}$.
\end{theorem}

\begin{proof}
The condensate Hilbert space $\mathcal{H}=L^2(\mathcal{C}_Q^{\rm red},d\mu)$
(Paper~X~\cite{Paper10}) carries a unitary representation of $2I$
induced by the $2I$-action on the reduced configuration space $\mathcal{C}_Q^{\rm red}$.
For $Q=2$, the configuration space $\mathcal{C}_{Q=2}^{\rm red}$ is the
space of $Q=2$ Hopf field configurations modulo translations.

The rotational degrees of freedom of the $Q=2$ gradient-knot
are parametrised by $SO(3)\cong\mathbb{RP}^3$.
The fundamental group $\pi_1(SO(3))=\mathbb{Z}_2$ implies
that the quantum states split into two sectors:
those with periodic boundary conditions around a $2\pi$ rotation loop
(integer spin) and those with anti-periodic boundary conditions
(half-integer spin).
For the $Q=2$ Hopfion, the natural geometric quantisation selects the
\emph{half-integer} sector because the Hopf fibration
$S^3\to S^2$ with fibre $S^1$ is itself a non-trivial $\mathbb{Z}_2$ fibration:
the parallel transport around a $2\pi$ rotation loop acquires holonomy $-1$
from the Hopf fibre phase $e^{i\pi Q/2}\big|_{Q=2}=e^{i\pi}=-1$.

Therefore the $Q=2$ gradient-knot transforms under the fundamental
spinor representation of $2I$, with the central element $z$ acting as $-1$
(equation~\eqref{P13:eq:z_action}).
Since $e^{i\cdot 2\pi J}=-1$ implies $J=\tfrac{1}{2}$ (mod 1),
the gradient-knot carries spin $J=\tfrac{1}{2}$.
\qed
\end{proof}

\begin{remark}[Physical origin of the $-1$ phase]
\label{P13:rem:hopf_phase}
The Hopf fibre phase $e^{i\pi Q/2}$ arises from the Wess--Zumino--Witten
term in the condensate effective action.
For the $Q=2$ configuration, a $2\pi$ rotation translates to a shift
of the WZW phase by $\pi Q/2 = \pi$, giving the holonomy $e^{i\pi}=-1$.
This is the analogue of the Finkelstein--Rubinstein mechanism for Skyrmions:
the quantisation of topological solitons in 3+1 dimensions selects
integer or half-integer spin according to the parity of the Hopf charge.
\end{remark}

\subsection{Fermi statistics and the Pauli exclusion principle}
\label{P13:sec:pauli}

\begin{corollary}[Fermi statistics]
\label{P13:cor:fermi_stats}
The two-particle wavefunction of two identical $Q=2$ gradient-knots is
antisymmetric under exchange:
\begin{equation}
  \boxed{\Psi(\mathbf{x}_2,\mathbf{x}_1) \;=\; -\Psi(\mathbf{x}_1,\mathbf{x}_2).}
  \label{P13:eq:antisymmetry}
\end{equation}
\end{corollary}

\begin{proof}
By the spin-statistics theorem in 3+1 dimensions, the exchange of two
identical particles with spin $J$ is topologically equivalent to a
$2\pi$ rotation of one particle, which introduces the phase $e^{i2\pi J}$.
From Theorem~\ref{P13:thm:spin_half}, $J=\tfrac{1}{2}$ and
$e^{i2\pi\cdot\frac{1}{2}}=e^{i\pi}=-1$.
Therefore the exchange of two $Q=2$ gradient-knots introduces a phase $-1$,
giving the antisymmetry relation~\eqref{P13:eq:antisymmetry}.
\qed
\end{proof}

\begin{corollary}[Pauli exclusion principle]
\label{P13:cor:pauli_exclusion}
No two identical $Q=2$ gradient-knots can occupy the same single-particle
quantum state.
\end{corollary}

\begin{proof}
Suppose two identical electrons occupy the same state, so that
$\mathbf{x}_1=\mathbf{x}_2$ (identical spatial and spin quantum numbers).
Substituting into~\eqref{P13:eq:antisymmetry}:
\begin{equation}
  \Psi(\mathbf{x},\mathbf{x}) = -\Psi(\mathbf{x},\mathbf{x})
  \quad\Longrightarrow\quad
  \Psi(\mathbf{x},\mathbf{x}) = 0.
  \label{P13:eq:pauli_zero}
\end{equation}
The two-particle amplitude vanishes identically whenever both particles
are in the same state.
\qed
\end{proof}

\begin{remark}[Common origin with the colour level shift $\delta n^{(C)}=\tfrac{1}{2}$]
\label{P13:rem:delta_n_half_origin}
Paper~V~\cite{Paper5}, Proposition~7.1 derives the
universal colour level shift $\delta n^{(C)}=\tfrac{1}{2}$ from the
$\mathrm{CP}^2$ fibre conformal weight $h=\tfrac{1}{2}$.
The $\tfrac{1}{2}$ there and the spin $J=\tfrac{1}{2}$ here share a
common algebraic origin: both arise because the central element
$z\in Z(2I)$ acts as $-1$ on half-integer representations of $2I$.

In the colour sector (Paper~V): the $\mathrm{CP}^2$ fibre of the Hopf
fibration $S^1\to S^5\to\mathrm{CP}^2$ transforms under the fundamental
spinor of $2I$, giving conformal weight shift
$\delta h = \tfrac{1}{2}$.

In the rotation sector (this section): the $Q=2$ gradient-knot transforms
under the same fundamental spinor of $2I\subset SU(2)$, giving spin
$J=\tfrac{1}{2}$ and hence Fermi statistics.

The electron is fermionic and quarks carry colour because the same
$\mathbb{Z}_2$ structure of $2I$ governs both sectors.
Both the Pauli exclusion principle of atomic physics and the colour
structure of hadronic physics have a single topological root in the
binary icosahedral symmetry of the condensate.
\end{remark}

\subsection{Complete spin assignment from the Q-tower}
\label{P13:sec:spin_tower}

Theorem~\ref{P13:thm:spin_half} derives $J=\tfrac{1}{2}$ for $Q=2$ by
minimising the conformal weight at the first non-trivial half-integer spin.
The same argument extends to all Hopf charges.

\begin{theorem}[Spin from Hopf charge: $J=Q/4$]
\label{P13:thm:spin_tower}
A gradient-knot with Hopf charge $Q$ in the condensate ground state
carries spin
\begin{equation}
  \boxed{J \;=\; \frac{Q}{4},}
  \label{P13:eq:J_Q4}
\end{equation}
transforming in the $(Q/2+1)$-dimensional irreducible representation
of $2I\subset SU(2)$.
\end{theorem}

\begin{proof}
\textbf{Step 1 (Berry phase).}
By the Hopf fibre holonomy (Theorem~\ref{P13:thm:spin_half}, extended to
general $Q$), the Berry phase under a $2\pi$ rotation of the $Q$-grade
gradient-knot is
\begin{equation}
  \theta_{\rm Berry}(2\pi) \;=\; e^{i\pi Q/2}.
  \label{P13:eq:berry_general}
\end{equation}
Since $e^{i\cdot 2\pi J} = \theta_{\rm Berry}(2\pi)$,
this constrains $J = Q/4$ modulo integers:
\begin{equation}
  J \;\in\; \Bigl\{\tfrac{Q}{4},\;\tfrac{Q}{4}+1,\;\tfrac{Q}{4}+2,\;\ldots\Bigr\}.
  \label{P13:eq:J_candidates}
\end{equation}

\textbf{Step 2 (Energy minimisation).}
In the condensate's WZW boundary theory at level $k=3$ (Paper~I~\cite{Paper1}),
the conformal weight of the spin-$J$ primary operator is
\begin{equation}
  h(J) \;=\; \frac{J(J+1)}{k+2} \;=\; \frac{J(J+1)}{5}.
  \label{P13:eq:conformal_weight}
\end{equation}
Since $h(J)$ is strictly increasing in $J$ for $J\ge0$, the ground state
at fixed $Q$ minimises $h$ subject to the constraint~\eqref{P13:eq:J_candidates}.
The minimum is achieved at the smallest allowed value:
$J_{\min} = Q/4$.

\textbf{Step 3 (Representation).}
At $J = Q/4$, the representation of $2I\subset SU(2)$ has dimension
$2J+1 = Q/2+1$.
This dimension is an irreducible representation of $2I$
for $Q = 0, 2, 4, 6, 8$, confirmed by the character table of $2I$
(Paper~V~\cite{Paper5}; see also Table~\ref{P13:tab:spin_assignments}).
\qed
\end{proof}

\begin{table}[h!]
\centering
\caption{Spin assignments from the $Q$-tower. The conformal weight
is $h = J(J+1)/5$ at WZW level $k=3$. The $2I$-irrep dimension is $2J+1$.
The Q=6 slot has no Standard Model assignment (it would be the gravitino
in supersymmetric extensions).}
\label{P13:tab:spin_assignments}
\begin{tabular}{cccccl}
\toprule
$Q$ & $J=Q/4$ & $2J+1$ & $h$ & Statistics & SM identification \\
\midrule
$0$ & $0$ & $1$ & $0$    & bosonic  & Higgs $H^0$ \\
$2$ & $\tfrac{1}{2}$ & $2$ & $\tfrac{3}{20}$ & fermionic & $e,\mu,\tau,\nu$; quarks \\
$4$ & $1$ & $3$ & $\tfrac{2}{5}$ & bosonic  & $\gamma,W^\pm,Z^0$; gluons \\
$6$ & $\tfrac{3}{2}$ & $4$ & $\tfrac{3}{4}$ & fermionic & --- (gravitino slot) \\
$8$ & $2$ & $5$ & $\tfrac{6}{5}$ & bosonic  & graviton \\
\bottomrule
\end{tabular}
\end{table}

\begin{remark}[The Higgs boson as the $Q=0$ condensate modulus]
\label{P13:rem:higgs_q0}
The condensate density $\rho(x)$ is a Lorentz scalar field.
Its amplitude excitation $\delta\rho(x)$ around the condensate VEV
(the Higgs boson) is therefore also a scalar: $J_{\rm Higgs}=0$.
In terms of the $Q$-tower, $Q=0$ corresponds to the trivial homotopy
class (no Hopf knot): the Higgs is not a topological soliton but rather
the radial fluctuation of the condensate field, consistent with
$e^{i\pi\cdot0/2} = +1$ (bosonic, $J=0$).
The condensate framework therefore derives the Higgs spin of zero rather
than postulating it.
\end{remark}

\begin{remark}[Gauge bosons as $Q=4$, $J=1$]
\label{P13:rem:gauge_q4}
Gauge bosons (photon, $W^\pm$, $Z^0$, gluons) are spin-1 connections
in the condensate's fibre bundles.
Their spin $J=1$ is consistent with $Q=4$ (the minimal non-trivial bosonic
Hopf charge), placing them in the $3$-dimensional ($J=1$) irreducible
representation of $I\subset SO(3)$.
The conformal weight $h=2/5$ at the $Q=4$ level is larger than the
fermionic $h=3/20$ at $Q=2$, consistent with gauge bosons being
``composite'' objects (binding two $Q=2$ knots mediated by the
condensate gauge connection).
The specific gauge group structure ($\mathrm{U}(1)$, $\mathrm{SU}(2)$,
$\mathrm{SU}(3)_C$) is established separately in Papers~I and~V;
here we derive only the universal spin-$1$ assignment.
\end{remark}

\begin{remark}[Graviton at $Q=8$, $J=2$]
\label{P13:rem:graviton_q8}
The graviton carries $J=2$, consistent with $Q=8$ and the
$5$-dimensional ($J=2$) representation of $I\subset SO(3)$.
Paper~VII~\cite{Paper7} derives gravity as emergent from the condensate
metric; whether the graviton is a fundamental $Q=8$ gradient-knot or
an acoustic mode of the condensate (both giving $J=2$) is not
resolved here.
The conformal weight $h=6/5$ at $Q=8$ exceeds unity, placing
the graviton above the unitarity bound of the $k=3$ WZW theory
(where $h_{\max}=J_{\max}(J_{\max}+1)/5$ with $J_{\max}=k/2=3/2$ for
ordinary WZW primaries at $k=3$).
This is consistent with the graviton being outside the renormalisable
sector of the condensate's boundary CFT, which encodes its
non-renormalisable nature in the conventional framework.
\end{remark}

\begin{remark}[The $Q=10$ condensate and its spin representation]
\label{P13:rem:condensate_q10}
The condensate itself has Hopf charge $Q=10$ (Paper~I~\cite{Paper1}).
By equation~\eqref{P13:eq:J_Q4}, this corresponds to $J=5/2$ and the
$6$-dimensional half-integer representation of $2I$, which is
the \emph{largest} irreducible representation appearing in
Table~\ref{P13:tab:spin_assignments} and the character table of $2I$.
This is consistent with the condensate being the ``maximal'' object
in the framework — all lower-$Q$ particles are excitations of it —
and with the $6$-dimensional rep of $2I$ exhausting the half-integer
sector's largest representation.
The condensate VEV $\langle\rho\rangle$ is of course a scalar
(Lorentz singlet); the $J=5/2$ assignment refers to the transformation
properties of the full $Q=10$ gradient-knot under the icosahedral
symmetry of the condensate, not to its Lorentz spin.
\end{remark}

\begin{remark}[Scope: fermionic statistics from the $Q$-tower]
\label{P13:rem:scope_fermi}
The $Q=2$ gradient-knot at tower level $n=20$ is the electron.
By Theorem~\ref{P13:thm:spin_tower}, any $Q=2$ gradient-knot at any
tower level $n$ carries $J=\tfrac{1}{2}$ and obeys Fermi statistics:
muons ($n=18$), tau leptons ($n=16$), and quarks with colour
(Paper~V~\cite{Paper5}) are therefore all fermionic.
For $Q=4$, Theorem~\ref{P13:thm:spin_tower} gives $J=1$ (bosonic);
for $Q=8$, $J=2$ (also bosonic, graviton).
The complete spin assignment is given in Table~\ref{P13:tab:spin_assignments}.
\end{remark}

%══════════════════════════════════════════════════════════════════════════════
\section{Open Problems}
\label{P13:sec:open}
%══════════════════════════════════════════════════════════════════════════════

All theoretical open problems identified at the start of this project
have been resolved in the course of the paper.
The single remaining item is experimental.

\begin{enumerate}[noitemsep]

\item \textbf{[Experimental] Penning-trap mass anisotropy measurement.}
The $2.13\%$ free-electron mass anisotropy
$\Delta m/m_e = 1/\vp^8$ (Proposition~\ref{P13:prop:mass_anisotropy})
is the only prediction of this paper accessible to near-term experiment.
A dedicated Penning-trap programme should:
(a) identify the local condensate gradient direction
(likely aligned with the CMB dipole or galactic centre);
(b) measure the cyclotron frequency $\omega_c = eB/m_{\rm eff}$ as a
function of trap orientation over a full annual cycle;
(c) search for a $2.13\%$ sinusoidal modulation.
Current $g$-factor measurements (Hanneke et al.~\cite{Hanneke2008})
achieve fractional precision $\sim10^{-13}$;
in principle they could detect a $2\%$ anisotropy if correctly modelled.

\end{enumerate}

\noindent
The resolved items are recorded here for completeness:
\begin{enumerate}[noitemsep,label=(\roman*)]
\item Schr\"odinger equation from condensate quantisation, with BO coupling
  $d_{0k}=0$ exactly and error $O(\eps_{\rm cond})\sim10^{-27}$
  (Theorem~\ref{P13:thm:schrodinger}, Theorem~\ref{P13:thm:NA_vanishing},
  Corollary~\ref{P13:cor:improved_BO}).
\item Condensate chameleon screening at the QCD scale
  (Paper~V~\cite{Paper5}, Proposition~6.9).
\item Spin-$\tfrac{1}{2}$, Fermi statistics, Pauli exclusion, and
  $J=Q/4$ for all SM particles
  (Theorems~\ref{P13:thm:spin_half}--\ref{P13:thm:spin_tower}).
\item Condensate corrections to multi-electron atoms ($Z>1$):
  spherical harmonics universal, $\Delta E\lesssim10^{-24}\,\mathrm{eV}$,
  Hartree--Fock exact
  (Propositions~\ref{P13:prop:multi_electron_screening},
  \ref{P13:prop:multi_electron_energy}).
\end{enumerate}

%══════════════════════════════════════════════════════════════════════════════
\section{Summary}
\label{P13:sec:summary}
%══════════════════════════════════════════════════════════════════════════════

This paper derives ten results constituting a complete foundation for
atomic and molecular quantum mechanics from the density-feedback
Faddeev--Niemi Hopfion condensate.
All results follow without additional postulates.

\paragraph{The condensate-modified Hamiltonian.}
The suppression law $\Seff(\theta,\rho)=\sin^4\!\theta/[\vp^6(1+\bst\rho)]$
adds a single correction to the hydrogen Hamiltonian:
\begin{equation*}
  \Heff
  = H_{\rm Coulomb}
  - \frac{\hbar^2}{2m_e}\,\frac{\epsr}{r^2}\,L^2,
  \qquad
  \epsr = \frac{1}{\vp^6(1+\bst\rho(r))},
\end{equation*}
acting only on angular kinetic energy (radial correction is exactly zero).

\paragraph{Chameleon screening and spherical harmonics.}
Inside atoms the Coulomb field drives
$\bst\rho_{\rm atom}(a_0)\approx2.6\times10^{25}$,
suppressing the correction to $\eps(a_0)\approx2.2\times10^{-27}$.
The effective Hamiltonian reduces to the standard Coulomb Hamiltonian,
giving spherical harmonics $Y_\ell^m$ as exact orbital shapes ---
derived from the chameleon mechanism, not assumed.
The condensate energy correction to hydrogen levels is $\sim10^{-24}\,\mathrm{eV}$
($\sim10^{20}$ below the fine structure); it breaks SO(4) degeneracy but is unobservable.

\paragraph{Universality for all atoms.}
The chameleon screening satisfies $\bst\rho_{\rm atom}^{(Z)}\ge\bst\rho_{\rm atom}^{(1)}\gg1$
throughout the electronic density of every $Z$-electron atom, because
$Z_{\rm eff}(r)\ge1$ always.
Spherical harmonics are therefore exact for all atoms and all $Z$.
The condensate energy correction scales as $Z^{-4}$ for inner electrons
and is bounded by $\sim10^{-24}\,\mathrm{eV}$ for valence electrons.
The Hartree--Fock equations are exact within the condensate framework.

\paragraph{The free-electron mass anisotropy.}
In the cosmic condensate background at $\bst\rhoinf=\vp$ (exact),
using $1+\vp=\vp^2$:
\[
  \frac{\Delta m}{m_e} = \frac{1}{\vp^8} \approx 2.13\%,
  \qquad
  \frac{m_{\rm tangential}}{m_{\rm radial}}
  = \frac{\vp^8+1}{\vp^8} \approx 1.02129,
\]
a falsifiable Penning-trap prediction.

\paragraph{The Schr\"odinger equation.}
The Schr\"odinger equation is derived from Paper~X's condensate quantum
theory via the translation fibration $\mathcal{C}_Q^{\rm red}=\mathbb{R}^3\times\mathcal{C}_{\rm internal}$,
the Noether momentum, Lorentz covariance fixing $M_{\rm eff}=m_e$,
and the Born--Oppenheimer separation at $\eps_{\rm BO}=\aem^2/(2\vp)\approx1.6\times10^{-5}$.

The formal BO bound $O(\eps_{\rm BO}^2)\approx2.7\times10^{-10}$
is improved dramatically: the derivative coupling $d_{0k}=\langle k_{\rm int}|\partial/\partial R|0_{\rm int}\rangle=0$
exactly for all $k$, because internal states are $R$-independent by
the translation fibration.
The actual error is $O(\eps_{\rm cond})\approx10^{-27}$,
a factor of $\sim10^{17}$ tighter than the formal bound.
The Schr\"odinger equation is derived, not postulated, and is exact to
one part in $10^{27}$.

\paragraph{Fermionic statistics, Pauli exclusion, and the spin tower.}
The binary icosahedral group $2I$ is a double cover of $I$;
its central element $z$ acts as $-1$ on the fundamental spinor representation.
The $Q=2$ gradient-knot transforms in this representation, giving
spin $J=\tfrac{1}{2}$ (Theorem~\ref{P13:thm:spin_half}), Fermi statistics
(Corollary~\ref{P13:cor:fermi_stats}), and Pauli exclusion
(Corollary~\ref{P13:cor:pauli_exclusion}).

The general formula $J=Q/4$ (Theorem~\ref{P13:thm:spin_tower}) extends
this to all Hopf charges, with the ground-state spin selected by
minimising the WZW conformal weight $h=J(J+1)/5$ at level $k=3$:
\begin{center}
$J=0$ (Higgs, $Q=0$),\quad
$J=\tfrac{1}{2}$ (fermions, $Q=2$),\quad
$J=1$ (gauge bosons, $Q=4$),\quad
$J=2$ (graviton, $Q=8$).
\end{center}
The same $\mathbb{Z}_2$ structure gives the colour level shift
$\delta n^{(C)}=\tfrac{1}{2}$ of Paper~V~\cite{Paper5}.

\paragraph{Unified mechanism.}
The single denominator $(1+\bst\rho)$ unifies all results:
it resolves the cosmological constant hierarchy at $\rho=\rhoinf$
(Paper~VII~\cite{Paper7}), forces orbital shapes to spherical harmonics
at $\rho=\rho_{\rm atom}$, produces the $1/\vp^8$ mass anisotropy at
$\rho=\rhoinf$, and underlies the $Q=2$ fermionic statistics via the
$2I$ double-cover structure.
All of standard atomic and molecular quantum mechanics follows from
the condensate without additional postulates.


\begin{thebibliography}{99}

\bibitem{Paper1}
F.~Manfredi,
\textit{The Density-Feedback Faddeev--Niemi Hopfion},
Zenodo (2026).
\href{https://doi.org/10.5281/zenodo.19342027}{doi:10.5281/zenodo.19342027}

\bibitem{Paper4}
F.~Manfredi,
\textit{The Hopf Spoke, WZW Fermion Mass Renormalization, and
  Three- and Four-Term Formulas for the Fine Structure Constant},
Zenodo (2026).
\href{https://doi.org/10.5281/zenodo.19504729}{doi:10.5281/zenodo.19504729}

\bibitem{Paper5}
F.~Manfredi,
\textit{Colour Confinement, the QCD Scale, and Hadronic Structure
  from the Density-Feedback Hopfion},
Zenodo (2026).
\href{https://doi.org/10.5281/zenodo.19638857}{doi:10.5281/zenodo.19638857}

\bibitem{Paper7}
F.~Manfredi,
\textit{Emergent Gravity, Inflation, Dark Energy, and the Weak
  Equivalence Principle from the Density-Feedback Hopfion Condensate},
Zenodo (2026).
\href{https://doi.org/10.5281/zenodo.19646953}{doi:10.5281/zenodo.19646953}

\bibitem{Paper10}
F.~Manfredi,
\textit{Quantisation of the Hopfion Condensate: Hilbert Space,
  Liouville Measure, and the Born Rule},
Zenodo (2026).
\href{https://doi.org/10.5281/zenodo.20075279}{doi:10.5281/zenodo.20075279}

\bibitem{Paper11}
F.~Manfredi,
\textit{The Golden-Spiral as a Universal Mass Hierarchy Map:
  From Condensate Scale to Chemical Periodicity},
Zenodo (2026).
\href{https://doi.org/10.5281/zenodo.20173763}{doi:10.5281/zenodo.20173763}

\bibitem{Paper12}
F.~Manfredi,
\textit{The Profile Normalisation Conjecture ---
  Proof via CMB Energy Identity and Two-Route Equivalence},
Zenodo (2026).
\href{https://doi.org/10.5281/zenodo.20210460}{doi:10.5281/zenodo.20210460}

\bibitem{Hanneke2008}
D.~Hanneke, S.~Fogwell, and G.~Gabrielse,
\textit{New Measurement of the Electron Magnetic Moment and
  the Fine Structure Constant},
Phys.\ Rev.\ Lett.\ \textbf{100}, 120801 (2008).
\href{https://doi.org/10.1103/PhysRevLett.100.120801}{doi:10.1103/PhysRevLett.100.120801}

\end{thebibliography}

\end{document}